% !TEX root = main.tex
\chapter{COMBINATORIAL STRUCTURES AND ORDINARY GENERATING FUNCTIONS}
\label{chap:combinatorial structures and ordinary generating functions}

\section{Symbolic enumeration methods}
% Note I.2
\noteblock{I.2}{Unimodal permutations.}

A unimodal permutation climbs to a single peak and then descends. In symbolic terms,
\begin{equation*}
    \mathcal{P}^{\mathcal{U}} = \mathrm{S}\textsc{et}(\mathcal{Z}) \star \mathcal{Z}^{\blacksquare} \star \mathrm{S}\textsc{et}(\mathcal{Z})
\end{equation*}
\begin{equation*}
    P^{U}(z) = \int_{0}^{z}\left(e^{t} \cdot \frac{\dd}{\dd t} t \cdot e^{t}\right)\dd t = \frac{1}{2}(e^{2z} - 1)
\end{equation*}
which immediately yields
\begin{equation*}
    P^{U}_{n} = n![z^{n}] \frac{1}{2} e^{2z} = n!\frac{1}{2} \frac{2^{n}}{n!} = 2^{n - 1}
\end{equation*}
For a direct counting argument,
\begin{equation*}
    P^{U}_{n} = \sum_{k = 1}^{n} \binom{n - 1}{k - 1} = \sum_{j = 0}^{n - 1} \binom{n - 1}{j} = 2^{n - 1}
\end{equation*}
where $k$ is the position of $n$ in the permutation.

\section{Admissible constructions and specifications}

\section{Integer compositions and partitions}
% Note I.16
\noteblock{I.16}{Partitions with summands bounded in number and size.}

Consider partitions of size $n$ with at most $k$ summands, each bounded by $l$. Their generating function is
\begin{equation} \label{eq:I.16.1}
    [z^{n}] \frac{(1 - z)(1 - z^{2}) \cdots (1 - z^{k + l})}{((1 - z)(1 - z^{2}) \cdots (1 - z^{k})) \cdot ((1 - z)(1 - z^{2}) \cdots (1 - z^{l}))}
\end{equation}
Let $f(k, l)$ denote this generating function, so the coefficient of $z^{n}$ counts the admissible partitions of $n$.
The natural split is by the largest part: it is either strictly smaller than $l$, or exactly $l$. Thus
\begin{equation} \label{eq:I.16.2}
    f(k, l) = f(k, l - 1) + z^{l}f(k - 1, l)
\end{equation}
To package the dependence on $l$, write
\begin{equation*}
    F_{k}(x) = \sum_{l = 0}^{\infty} f(k, l)x^{l}
\end{equation*}
For the base case $k = 0$, only the empty partition survives, so $f(0,l) = 1$ for all $l \geq 0$. Hence
\begin{equation*}
    F_{0}(x) = \sum_{l = 0}^{\infty}1 \cdot x^{l} = \frac{1}{1 - x}
\end{equation*}
Now multiply \ref{eq:I.16.2} by $x^{l}$ and sum over $l \geq 1$:
\begin{equation*}
    \sum_{l = 1}^{\infty} f(k, l)x^{l} - \sum_{l = 1}^{\infty} f(k, l - 1)x^{l} = \sum_{l = 1}^{\infty} f(k - 1, l)(zx)^{l}
\end{equation*}
\begin{itemize}
    \item The first term simplifies to $F_{k}(x) - 1$, since $f(k, 0) = 1$.
    \item The second term is simply $x F_{k}(x)$.
    \item The right-hand side yields $F_{k - 1}(zx) - 1$.
\end{itemize}
Substituting these pieces back in gives
\begin{equation*}
    (1 - x)F_{k}(x) - 1 = F_{k - 1}(zx) - 1 \implies F_{k}(x) = \frac{1}{1 - x} F_{k - 1}(zx)
\end{equation*}
Iterating the relation gives
\begin{equation*}
    \begin{aligned}
        F_{k}(x)
        &= \frac{1}{1 - x} F_{k - 1}(zx) \\
        &= \frac{1}{1 - x} \frac{1}{1 - zx} F_{k - 2}(z^{2}x) \\
        &= \frac{1}{(1 - x)(1 - zx) \cdots (1 - z^{k - 1}x)}F_{0}(z^{k}x)
    \end{aligned}
\end{equation*}
Since $F_{0}(x) = \frac{1}{1 - x}$, we have $F_{0}(z^{k}x) = \frac{1}{1 - z^{k}x}$. Thus
\begin{equation*}
    F_{k}(x) = \sum_{l = 0}^{\infty} f(k, l)x^{l} = \frac{1}{(1 - x)(1 - zx) \cdots (1 - z^{k}x)}
\end{equation*}
Since the denominator has distinct linear factors, we may use partial fractions:
\begin{equation*}
    \frac{1}{\prod_{i = 0}^{k}(1 - z^{i}x)} = \sum_{j = 0}^{k} \frac{A_{j}}{1 - z^{j}x}
\end{equation*}
To isolate $A_{j}$, multiply through by $(1 - z^{j}x)$ and then set $x = z^{-j}$:
\begin{equation*}
    \begin{aligned}
        A_{j}
        &= \left[\frac{1 - z^{j}x}{\prod_{i = 0}^{k}(1 - z^{i}x)}\right]_{x = z^{-j}} \\
        &= \frac{1}{\prod_{i = 0}^{j - 1}(1 - z^{i}z^{-j})\prod_{i = j + 1}^{k}(1 - z^{i}z^{-j})} \\
        &= \frac{(-1)^{j}z^{j(j + 1)/2}}{\prod_{i = 1}^{j}(1 - z^{i})\prod_{i = 1}^{k - j}(1 - z^{i})}
    \end{aligned}
\end{equation*}
The coefficient of $x^{l}$ in $F_{k}(x)$ is then read off term by term. For $\frac{A_{j}}{1 - z^{j}x}$, it is simply $A_{j}(z^{j})^{l}$.
\begin{equation*}
    \begin{aligned}
        f(k, l)
        &= \sum_{j = 0}^{k} A_{j}(z^{j})^{l} \\
        &= \sum_{j = 0}^{k} \frac{(-1)^{j}z^{j(j + 1)/2} \cdot z^{jl}}{\prod_{i = 1}^{j}(1 - z^{i})\prod_{i = 1}^{k - j}(1 - z^{i})} \\
        &= \frac{1}{\prod_{i = 1}^{k}(1 - z^{i})} \sum_{j = 0}^{k} \binom{k}{j}_{z}(-1)^{j}z^{j(j + 1)/2 + jl}
    \end{aligned}
\end{equation*}
where $\binom{k}{j}_{z}$ is the Gaussian binomial coefficient:
\begin{equation*}
    \begin{aligned}
        \binom{k}{j}_{z}
        &= \frac{\prod_{i = 1}^{k}(1 - z^{i})}{\prod_{i = 1}^{j}(1 - z^{i})\prod_{i = 1}^{k - j}(1 - z^{i})} \\
        &= \frac{(1 - z^{k})(1 - z^{k - 1}) \cdots (1 - z^{k - j + 1})}{(1 - z)(1 - z^{2}) \cdots (1 - z^{j})}
    \end{aligned}
\end{equation*}
A small rearrangement of the exponent gives
\begin{equation*}
    \frac{j(j + 1)}{2} + jl = \frac{j(j - 1)}{2} + j(l + 1)
\end{equation*}
so the sum becomes
\begin{equation*}
    f(k, l) = \frac{1}{\prod_{i = 1}^{k}(1 - z^{i})} \sum_{j = 0}^{k} \binom{k}{j}_{z}(-1)^{j}z^{j(j - 1)/2}(z^{l + 1})^{j}
\end{equation*}
The $q$-binomial theorem states that for a variable $X$:
\begin{equation} \label{eq:I.16.3}
    \prod_{i = 0}^{k - 1}(1 - z^{i}X) = \sum_{j = 0}^{k} \binom{k}{j}_{z}(-1)^{j}z^{j(j - 1)/2}X^{j}
\end{equation}
Substituting $X = z^{l + 1}$ into the theorem recovers \ref{eq:I.16.1}:
\begin{equation*}
    \begin{aligned}
        f(k, l)
        &= \frac{(1 - z^{l + 1})(1 - z^{l + 2}) \cdots (1 - z^{l + k})}{(1 - z^{1})(1 - z^{2}) \cdots (1 - z^{k})} \\
        &= \frac{(1 - z)(1 - z^{2}) \cdots (1 - z^{k + l})}{((1 - z)(1 - z^{2}) \cdots (1 - z^{k})) \cdot ((1 - z)(1 - z^{2}) \cdots (1 - z^{l}))}
    \end{aligned}
\end{equation*}
For completeness, let us also justify \ref{eq:I.16.3}. Set
\begin{equation*}
    P_{k}(X) = \prod_{i = 0}^{k - 1}(1 - z^{i}X)
\end{equation*}
We claim that $P_{k}(X) = \sum_{j = 0}^{k} \binom{k}{j}_{z}(-1)^{j}z^{j(j - 1)/2}X^{j}$.
Start from the elementary relation
\begin{equation*}
    P_{k}(X) = (1 - z^{k - 1}X)P_{k - 1}(X)
\end{equation*}
Next, replace $X$ by $zX$:
\begin{equation*}
    P_{k}(zX) = \prod_{i = 0}^{k - 1}(1 - z^{i + 1}X) = \frac{1 - z^{k}X}{1 - X}P_{k}(X)
\end{equation*}
After rearranging, this becomes the functional equation
\begin{equation*}
    (1 - X)P_{k}(zX) = (1 - z^{k}X)P_{k}(X)
\end{equation*}
Let $P_{k}(X) = \sum_{j = 0}^{k} c_{j}X^{j}$. We substitute this power series into the functional equation:
\begin{equation*}
    (1 - X)\sum_{j = 0}^{k} c_{j}(zX)^{j} = (1 - z^{k}X)\sum_{j = 0}^{k} c_{j}X^{j}
\end{equation*}
Comparing coefficients of $X^{j}$ on both sides yields:
\begin{equation*}
    c_{j}z^{j} - c_{j - 1}z^{j - 1} = c_{j} - c_{j - 1}z^{k}
\end{equation*}
so
\begin{equation*}
    c_{j} = c_{j - 1} \frac{z^{j - 1} - z^{k}}{z^{j} - 1} = -c_{j - 1}z^{j - 1} \frac{1 - z^{k - j + 1}}{1 - z^{j}}
\end{equation*}
Starting from $c_{0} = P_{k}(0) = 1$, we iterate:
\begin{equation*}
    c_{1} = -c_{0}z^{0} \frac{1 - z^{k}}{1 - z^{1}}, \quad
    c_{2} = -c_{1}z^{1} \frac{1 - z^{k - 1}}{1 - z^{2}} = (-1)^{2}z^{0 + 1} \frac{(1 - z^{k})(1 - z^{k - 1})}{(1 - z^{1})(1 - z^{2})}
\end{equation*}
The exponent of $z$ is $0 + 1 + 2 + \cdots + (j - 1) = \frac{j(j - 1)}{2}$, and the remaining fraction is exactly the Gaussian binomial coefficient. Therefore,
\begin{equation*}
    c_{j} = (-1)^{j}z^{j(j - 1)/2} \binom{k}{j}_{z}
\end{equation*}
Plugging the coefficients back into the expansion, we obtain
\begin{equation*}
    \prod_{i = 0}^{k - 1}(1 - z^{i}X) = \sum_{j = 0}^{k} \binom{k}{j}_{z}(-1)^{j}z^{j(j - 1)/2}X^{j}
\end{equation*}
which is exactly what we wanted.

% Note I.17
\noteblock{I.17}{Systems of linear Diophantine inequalities.}

This note contains five separate counting problems. They are related by theme, but each one has its own constraint set and should be read independently.

\subhead{Problem 1}
\begin{equation*}
    \left\{x_{1}\geq 0,\; x_{2}\geq 2x_{1},\; x_{3}\geq 2x_{2},\; x_{4}\geq 2x_{3}\right\}
\end{equation*}
At first glance this system looks mildly unfriendly, but slack variables settle it almost at once. Introduce non-negative variables
$y_{j} \in \mathbb{Z}_{\geq 0}$ to measure the ``excess'' of each term above its lower bound.
\begin{align*}
    x_{1} &= y_{1} \\
    x_{2} &= 2x_{1} + y_{2} = 2y_{1} + y_{2} \\
    x_{3} &= 2x_{2} + y_{3} = 4y_{1} + 2y_{2} + y_{3} \\
    x_{4} &= 2x_{3} + y_{4} = 8y_{1} + 4y_{2} + 2y_{3} + y_{4}
\end{align*}
The total weight is:
\begin{equation*}
    S = x_{1} + x_{2} + x_{3} + x_{4} = 15y_{1} + 7y_{2} + 3y_{3} + y_{4}.
\end{equation*}
Now the generating function is immediate: the variables $y_{j}$ are independent and each ranges freely over
the non-negative integers, so the OGF $F(z)$ is just the product of the four one-variable OGFs:
\begin{itemize}
    \item For $15y_{1}: \sum_{y_{1} = 0}^{\infty}(z^{15})^{y_{1}} = \frac{1}{1 - z^{15}}$
    \item For $ 7y_{2}: \sum_{y_{2} = 0}^{\infty}(z^{ 7})^{y_{2}} = \frac{1}{1 - z^{ 7}}$
    \item For $ 3y_{3}: \sum_{y_{3} = 0}^{\infty}(z^{ 3})^{y_{3}} = \frac{1}{1 - z^{ 3}}$
    \item For $  y_{4}: \sum_{y_{4} = 0}^{\infty}(z^{ 1})^{y_{4}} = \frac{1}{1 - z^{ 1}}$
\end{itemize}
Multiplying them together gives the answer:
\begin{equation*}
    F(z) = \frac{1}{(1 -z)(1 -z^{3})(1 -z^{7})(1 -z^{15})}
\end{equation*}

\partsep
\subhead{Problem 2}
\begin{equation*}
    \left\{x_{1} + x_{2}\leq x_{3}\right\}
\end{equation*}
This is a new problem, so we start again from its own constraint. Here slack variables turn the inequality into an equality with independent parts.
The constraint $x_{3}\geq x_{1} + x_{2}$ can be rewritten by introducing a non-negative integer $y\geq 0$:
\begin{equation*}
    x_{3} = x_{1} + x_{2} + y
\end{equation*}
The size of the composition is $S = x_{1} + x_{2} + x_{3} + x_{4} = 2x_{1} + 2x_{2} + y + x_{4}$.
Since the variables $x_{1}, x_{2}, y, x_{4}$ are independent and each ranges freely over the non-negative integers, the OGF is the product of
the generating functions for the individual terms in the sum:
\begin{itemize}
    \item For $2x_{1}: \frac{1}{1 - z^{2}}$
    \item For $2x_{2}: \frac{1}{1 - z^{2}}$
    \item For $y: \frac{1}{1 - z}$
    \item For $x_{4}: \frac{1}{1 - z}$
\end{itemize}
Multiplying these factors gives
\begin{equation*}
    F(z) = \frac{1}{(1 -z^{2})^{2}(1 -z)^{2}}
\end{equation*}

\partsep
\subhead{Problem 3}
\begin{equation*}
    \left\{x_{1} + x_{2}\geq x_{3}\right\}
\end{equation*}
This third problem is independent of the previous two. Since the inequality does not suggest a single slack variable, the cleanest route is to work with the complement.
Let $U(z) = (1 - z)^{-4}$ be the OGF for four unrestricted variables. The universe of compositions partitions into
those where $x_{1} + x_{2} < x_{3}$ and our target class $x_{1} + x_{2} \geq x_{3}$. We find the OGF for the complement
$x_{1} + x_{2} < x_{3}$. Set $x_{3} = x_{1} + x_{2} + 1 + y$, where $y\geq 0$.
The total sum $S$ is:
\begin{equation*}
    S = x_{1} + x_{2} + (x_{1} + x_{2} + 1 + y) + x_{4} = 2x_{1} + 2x_{2} + y + x_{4} + 1
\end{equation*}
The OGF for this complement, denoted by $C(z)$, is
\begin{equation*}
    C(z) = \frac{z}{(1 - z^{2})^{2}(1 - z)^{2}}
\end{equation*}
The OGF for our constraint is then
\begin{equation*}
    F(z) = U(z) - C(z) = \frac{1}{(1 - z)^{4}} - \frac{z}{(1 - z^{2})^{2}(1 - z)^{2}} = \frac{1 + z + z^{2}}{(1 - z)^{4}(1 + z)^{2}}
\end{equation*}

\partsep
\subhead{Problem 4}
\begin{equation*}
    \left\{x_{1} + x_{2}\leq x_{3} + x_{4}\right\}
\end{equation*}
We now move to a fourth, separate problem.
To find the OGF for $\left\{x_{1} + x_{2}\leq x_{3} + x_{4}\right\}$, combine the complement method with the symmetry of the variables. Let $A = x_{1} + x_{2}$ and $B = x_{3} + x_{4}$.
We seek the OGF for compositions satisfying $A\leq B$. In the universe of all compositions
$(x_{1}, x_{2}, x_{3}, x_{4})$, there are three mutually exclusive cases
$1.\; A < B,\; 2.\; A > B,\; 3.\; A = B$. By symmetry, the OGF for case 1 must be identical to the OGF for case 2,
since $A$ and $B$ are built from the same number of variables with the same weights.
Denote this common OGF by $G_{<}(z)$, and let $G_{=}(z)$ be the OGF for case 3.
The total unrestricted OGF is $U(z) = \frac{1}{(1 - z)^{4}}$. Therefore,
\begin{equation*}
    U(z) = 2G_{<}(z) + G_{=}(z)
\end{equation*}
Our target OGF is $F(z) = G_{<}(z) + G_{=}(z)$. Solving the system gives:
\begin{equation*}
    F(z) = \frac{U(z) + G_{=}(z)}{2}
\end{equation*}
Now $G_{=}(z)$ is the OGF for compositions in which the sum of the first two variables equals
the sum of the last two.
\begin{itemize}
    \item The number of compositions with $x_{1} + x_{2} = n$ is:
        $[z^{n}] \frac{1}{(1 - z)^{2}} = n + 1$.
    \item The number of compositions with $x_{3} + x_{4} = n$ is the same.
    \item For the equality $x_{1} + x_{2} = x_{3} + x_{4}$ to hold, both sums must equal some $n\geq 0$.
        \begin{equation*}
            G_{=}(z) = \sum_{n = 0}^{\infty}(n + 1)^{2}z^{2n} = \frac{1 + z^{2}}{(1 - z^{2})^{3}}
        \end{equation*}
\end{itemize}
Substituting $U(z)$ and $G_{=}(z)$ into our formula yields:
\begin{equation*}
    F(z) = \frac{1}{2} \left[\frac{1}{(1 - z)^{4}} + \frac{1 + z^{2}}{(1 - z^{2})^{3}}\right] = \frac{1 + z + 2z^{2}}{(1 - z)^{4}(1 + z)^{3}}
\end{equation*}

\partsep
\subhead{Problem 5}
\begin{equation*}
    \left\{x_{1}\leq x_{2},\; x_{2}\geq x_{3},\; x_{3}\leq x_{4}\right\}
\end{equation*}
The fifth problem is again independent.
Because the inequalities alternate direction, we cannot use the simple slack-variable chain that works for $x_{1}\leq x_{2}\leq x_{3}\leq x_{4}$. Instead, we use inclusion--exclusion.
We begin with the universe of all unrestricted compositions, whose OGF is $U(z) = \frac{1}{(1 - z)^{4}}$.
We wish to enforce three conditions
\begin{itemize}
    \item $C_{1}: x_{1}\leq x_{2}$
    \item $C_{2}: x_{2}\geq x_{3}$
    \item $C_{3}: x_{3}\leq x_{4}$
\end{itemize}
It is easier to work with chains of ``less than or equal'' inequalities, so consider the complement of the middle downward step.
Take the set of all compositions satisfying both $x_{1}\leq x_{2}$ and $x_{3}\leq x_{4}$.
From this set, we subtract those compositions that violate the middle constraint (i.e., $x_{2} < x_{3}$). \\
The OGF for the set defined by $x_{1}\leq x_{2}$ and $x_{3}\leq x_{4}$ is easy to obtain because the two pairs are independent.
\begin{itemize}
    \item For $x_{1}\leq x_{2}$, set $x_{2} = x_{1} + y_{1}$. The total contribution is $2x_{1} + y_{1}$,
        so the OGF is $\frac{1}{(1 - z^{2})(1 - z)}$.
    \item For $x_{3}\leq x_{4}$, the OGF is identically $\frac{1}{(1 - z^{2})(1 - z)}$.
\end{itemize}
Since the two pairs are independent, we multiply their OGFs:
\begin{equation*}
    A(z) = \frac{1}{(1 - z^{2})^{2}(1 - z)^{2}}
\end{equation*}
Now we find the OGF for the violation $x_{1}\leq x_{2} < x_{3}\leq x_{4}$ by introducing slack variables
\begin{itemize}
    \item $x_{1} = y_{1}$
    \item $x_{2} = y_{1} + y_{2}$
    \item $x_{3} = y_{1} + y_{2} + 1 + y_{3}$
    \item $x_{4} = y_{1} + y_{2} + 1 + y_{3} + y_{4}$
\end{itemize}
Summing the four expressions gives:
\begin{equation*}
    S = 4y_{1} + 3y_{2} + 2y_{3} + y_{4} + 2
\end{equation*}
The OGF for this violation is therefore:
\begin{equation*}
    V(z) = \frac{z^{2}}{(1 - z^{4})(1 - z^{3})(1 - z^{2})(1 - z)}
\end{equation*}
The OGF for the zigzag constraint is the OGF for the two independent pairs minus the OGF for the cases where the peak at $x_{2}$ fails:
\begin{equation*}
    \begin{aligned}
        F(z)
        &= \frac{1}{(1 - z^{2})^{2}(1 - z)^{2}} - \frac{z^{2}}{(1 - z^{4})(1 - z^{3})(1 - z^{2})(1 - z)} \\
        &= \frac{1}{(1 - z)(1 - z^{2})}\left[\frac{1}{(1 - z)(1 - z^{2})} - \frac{z^{2}}{(1 - z^{3})(1 - z^{4})}\right]
    \end{aligned}
\end{equation*}

% Note I.18
\noteblock{I.18}{Odd versus distinct summands.}

This bijection looks clever on first meeting and then almost inevitable in retrospect. We build a one-to-one correspondence between partitions into odd parts and partitions into distinct parts.
\begin{itemize}
    \item Odd $\rightarrow$ Distinct
        \begin{enumerate}
            \item Group identical odd parts together.
            \item If an odd part $k$ appears $m$ times, write $m$ in binary
                    $m = 2^{a_{1}} + 2^{a_{2}} + \cdots$ with distinct powers of two.
            \item Replace the $m$ copies of $k$ by $2^{a_{1}} \cdot k,\;
                    2^{a_{2}} \cdot k,\; \dots$
        \end{enumerate}
    \item Distinct $\rightarrow$ Odd
        \begin{enumerate}
            \item Factor each part as $2^{a}\times \text{odd}$.
            \item Split it into $2^{a}$ copies of that odd factor.
            \item Do this for all parts; now all parts are odd.
            \item Sort the resulting parts to obtain a partition into odd parts.
        \end{enumerate}
\end{itemize}
Here is a concrete pass through the map:
\begin{equation*}
    \begin{aligned}
        3 + 3 + 3 + 3 + 3 + 1 + 1
        &= 5\times 3 + 2\times 1 \\
        &= (2^{2} + 2^{0})\times 3 + 2^{1}\times 1 \\
        &= 4\times 3 + 1\times 3 + 2\times 1 \\
        &= 12 + 3 + 2
    \end{aligned}
\end{equation*}
\begin{equation*}
    \begin{aligned}
        12 + 3 + 2
        &= 4\times 3 + 3 + 2\times 1 \\
        &= (3 + 3 + 3 + 3) + 3 + (1 + 1) \\
        &= 3 + 3 + 3 + 3 + 3 + 1 + 1
    \end{aligned}
\end{equation*}

% Note I.19
\noteblock{I.19}{Euler’s pentagonal number theorem.}

To get a usable recurrence for the partition numbers, it is natural to look at the reciprocal of $P(z)$:
\begin{equation*}
    A(z) = \prod_{n = 1}^{\infty}(1 - z^{n})
\end{equation*}
Euler's Pentagonal Number Theorem says that this product collapses to a remarkably sparse series:
\begin{equation*}
    \begin{aligned}
        A(z)
            &= \sum_{k = -\infty}^{\infty}(-1)^{k}z^{k(3k + 1)/2} \\
            &= \sum_{k = -\infty}^{\infty}(-1)^{k}z^{k(3k - 1)/2} \\
            &= 1 - z^{1} - z^{2} + z^{5} + z^{7} - z^{12} - z^{15} + \cdots
    \end{aligned}
\end{equation*}
The exponents $g_{k} = \frac{k(3k - 1)}{2}$ are the generalized pentagonal numbers. Since $A(z)$ is the reciprocal of $P(z)$, we have
\begin{equation*}
    \sum_{n = 0}^{\infty} p(n)z^{n} \cdot \sum_{k = -\infty}^{\infty}(-1)^{k}z^{g_{k}} = 1
\end{equation*}
Comparing coefficients of $z^{n}$ now gives the recurrence that is genuinely useful in computation:
\begin{equation*}
    p(n) = p(n - 1) + p(n - 2) - p(n - 5) - p(n - 7) + p(n - 12) + p(n - 15) \cdots
\end{equation*}
A straightforward Python implementation is:
\newpage
\begin{lstlisting}
    def compute_partitions(N):
    P = [0] * (N + 1)
    P[0] = 1  # Base case: p(0) = 1

    for i in range(1, N + 1):
        k = 1
        while True:
            # Calculate pentagonal numbers for k and -k
            g_k_pos = k * (3 * k - 1) // 2
            g_k_neg = k * (3 * k + 1) // 2

            # Determine sign: +1 for k=1,2 (terms 1,2), -1 for k=3,4 (terms 3,4)...
            # Simplest way: (-1)^(k-1)
            sign = 1 if k % 2 == 1 else -1

            # Add p(i - g_k_pos)
            if g_k_pos <= i:
                P[i] += sign * P[i - g_k_pos]
            else:
                break # No more terms will fit

            # Add p(i - g_k_neg)
            if g_k_neg <= i:
                P[i] += sign * P[i - g_k_neg]
            else:
                break

            k += 1
    return P
\end{lstlisting}

% Note I.20
\noteblock{I.20}{A digital surprise.}
\begin{equation*}
    \begin{aligned}
        \varphi
            &= (1 - 10^{-1})(1 - 10^{-2})(1 - 10^{-3})(1 - 10^{-4}) \cdots \\
            &= \left.\prod_{n = 1}^{\infty}(1 - z^{n})\right|_{z = .1} \\
            &= \left.1 - z^{1} - z^{2} + z^{5} + z^{7} - z^{12} - z^{15} + z^{22} + z^{26} - \cdots \right|_{z = .1} \\
            &= .\dot{9} - .1 - .01 + .00001 + .0000001 - .000000000001 - \cdots
    \end{aligned}
\end{equation*}
The partial sums are
\begin{itemize}
    \item $\varphi_{1} = .99999 99999 99999 99999 99999$
    \item $\varphi_{2} = .89999 99999 99999 99999 99999$
    \item $\varphi_{3} = .88999 99999 99999 99999 99999$
    \item $\varphi_{4} = .89000 99999 99999 99999 99999$
    \item $\varphi_{5} = .89001 00999 99999 99999 99999$
    \item $\varphi_{6} = .89001 00999 98999 99999 99999$
    \item $\varphi_{7} = .89001 00999 98998 99999 99999$
    \item $\varphi_{8} = .89001 00999 98999 00000 00999$
\end{itemize}

% Note I.21
\noteblock{I.21}{Lattice points.}

We want the number of integer solutions to $\sum_{i = 1}^{d} x_{i}^{2} \le n^{2}$.
The natural object here is the theta series. For one coordinate $x_{i} \in \mathbb{Z}$, the contribution to the sum of squares is $x_{i}^{2}$, so the one-variable generating function is
\begin{equation*}
    \Theta(z) = \sum_{x \in \mathbb{Z}}z^{x^{2}} = 1 + 2 \sum_{n = 1}^{\infty} z^{n^{2}}
\end{equation*}
Because the coordinates are independent, the generating function for $\sum x_{i}^{2}$ is simply the $d$th power:
\begin{equation*}
    S(z) = (\Theta(z))^{d} = \sum_{k = 0}^{\infty} r_{d}(k)z^{k}
\end{equation*}
Here $[z^{k}] S(z)$, denoted $r_{d}(k)$, counts representations of $k$ as a sum of $d$ squares. Geometrically, that is the number of lattice points on the sphere of radius $\sqrt{k}$.
Therefore, the number of points $N(n, d)$ in the closed ball of radius $n$ is:
\begin{equation*}
    N(n, d) = \sum_{k = 0}^{n^{2}}r_{d}(k)
\end{equation*}
Passing from exact radius to radius at most $n$ is just a prefix sum, so we multiply by $\frac{1}{1 - z}$:
\begin{equation*}
    \sum_{k = 0}^{\infty} \left(\sum_{j = 0}^{k} r_{d}(j)\right)z^{k} = \frac{1}{1 - z} S(z) = \frac{1}{1 - z}(\Theta(z))^{d}
\end{equation*}
Extracting the coefficient of $z^{n^{2}}$ finishes the job:
\begin{equation*}
    N(n, d) = [z^{n^{2}}] \frac{1}{1 - z}(\Theta(z))^{d}
\end{equation*}

\section{Words and regular languages}
% Note I.23
\noteblock{I.23}{Alice, Bob, and coding bounds.}

By Proposition I.4, the class $\mathcal{C}$ of binary strings avoiding the pattern $11$ has generating function
\begin{equation*}
    C(z) = \frac{1 + z}{1 - z - z^{2}}
\end{equation*}
If codewords are allowed to have any length up to $L$, we need the sum of the first $L + 1$ coefficients. Its generating function is:
\begin{equation*}
    S(z) = \frac{C(z)}{1 - z} = \frac{1 + z}{(1 - z)(1 - z - z^{2})}
\end{equation*}
As usual, the nearest singularity controls the asymptotics. Since the pole at $z = 1/\varphi$ (with $\varphi = \frac{1+\sqrt{5}}{2}$) lies closer to the origin than the pole at $z = 1$, we obtain, as $L \to \infty$,
\begin{equation*}
    \sum_{j = 0}^{L} C_{j} \sim A \cdot \varphi^{L}
\end{equation*}
for some constant $A$. To encode $n$ bits of information, we want $2^{n} \leq \sum_{j = 0}^{L} C_{j}$, so taking base-2 logarithms gives:
\begin{equation*}
    n \leq \log_{2}\left(\sum_{j = 0}^{L} C_{j}\right) = L \log_{2}\varphi + O(1)
\end{equation*}
Rearranging for $L$, we obtain the lower bound:
\begin{equation*}
    L \geq \frac{n}{\log_{2}\varphi} + O(1) \approx 1.440420n + O(1)
\end{equation*}

% Note I.27
\noteblock{I.27}{Congruence languages.}

The DFA $\mathcal{A}$ for binary strings congruent to $2 \pmod{7}$ is the standard remainder automaton:
\begin{itemize}
    \item \textbf{States}: A set $\mathcal{Q} = \{S_{0}, \dots, S_{6}\}$, where each state $S_{r}$ represents the current remainder $r$ modulo 7.
    \item \textbf{Initial State}: $S_{0}$ (representing the value 0).
    \item \textbf{Accepting State}: $\overline{Q} = \{S_{2}\}$.
    \item \textbf{Alphabet}: $\mathcal{A} = \{0, 1\}$.
\end{itemize}
Each transition follows the familiar update rule $q_{r} \circ b = q_{(2r + b) \pmod{7}}$.
\begin{table}[H]
    \centering
    \caption{State Transition}
    \label{tab:flexible}
    \begin{tabularx}{\textwidth}{XXX}  % All center-aligned
        \toprule
        Current State (Remainder) & Input 0 ($2r \pmod{7}$) & Input 1 ($2r + 1 \pmod{7}$) \\
        \midrule
        $S_{0}$ & $S_{0}$ & $S_{1}$ \\
        $S_{1}$ & $S_{2}$ & $S_{3}$ \\
        $S_{2}$ & $S_{4}$ & $S_{5}$ \\
        $S_{3}$ & $S_{6}$ & $S_{0}$ \\
        $S_{4}$ & $S_{1}$ & $S_{2}$ \\
        $S_{5}$ & $S_{3}$ & $S_{4}$ \\
        $S_{6}$ & $S_{5}$ & $S_{6}$ \\
        \bottomrule
    \end{tabularx}
\end{table}

% Note I.29
\noteblock{I.29}{The forward equations.}

A last-letter decomposition gives the usual recurrence for the set $\mathcal{M}_{k}$ of words leading from the initial state $q_{0}$ to state $q_{k}$:
\begin{equation*}
    \mathcal{M}_{k} \cong \llbracket k = 0 \rrbracket \cdot \{\epsilon\} +
        \sum_{j, \alpha: q_{j} \circ \alpha = q_{k}} (\mathcal{M}_{j} \times \{\alpha\})
\end{equation*}
where $\llbracket \cdot \rrbracket$ is the Iverson bracket. Translating directly to generating functions,
\begin{equation*}
    M_{k}(z) = \delta_{k, 0} + z \sum_{j} M_{j}(z) T_{j, k}
\end{equation*}
In vector form, with $\mathbf{M}(z)$ as a column vector and $\mathbf{e}_{0}$ the first basis vector, this is simply the transposed analogue of the backward equations:
\begin{equation*}
    \mathbf{M}(z)^{T} = \mathbf{e}_{0}^{T} + z\mathbf{M}(z)^{T}T
\end{equation*}

% Note I.33.
\noteblock{I.33}{Waiting times in strings.}
Let us test the formulas on a deliberately tiny prefix language, $\hat{\mathcal{L}} = \{aa, bbb\}$.
\begin{enumerate}
    \item Prefix property: $aa$ is not a prefix of $bbb$, and conversely, so the language is prefix-free.
    \item Generating function: the OGF is
        $\hat{L}(z) = z^{2} + z^{3}$.
    \item Probability: $\hat{L}(1/2)$ is the probability that a random infinite string starts with a word of $\hat{\mathcal{L}}$:
        $\hat{L}(1/2) = (1/2)^{2} + (1/2)^{3} = 1/4 + 1/8 = 3/8$.
    \item Expected time: $\frac{1}{2} \hat{L}'(1/2)$ gives the expected matched length:
        $\hat{L}'(z) = 2z + 3z^{2} \implies \hat{L}'(1/2) = 1.75 \implies E = 0.875$.
    \item Unpacking the value:
        \begin{itemize}
            \item Word ``aa'': Probability $1/4$ of length 2 $\rightarrow$ Contribution: $2 \times 0.25 = 0.5$.
            \item Word ``bbb'': Probability $1/8$ of length 3 $\rightarrow$ Contribution: $3 \times 0.125 = 0.375$.
            \item No Match: Probability $5/8$ of length 0 $\rightarrow$ Contribution: $0 \times 0.625 = 0$.
        \end{itemize}
\end{enumerate}

% Note I.34.
\noteblock{I.34}{A probabilistic paradox on strings.}

Fix a binary pattern of length $k$. The generating function $T(z)$ for its first occurrence is:
\begin{equation*}
    T(z) = \frac{z^{k}}{z^{k} + (1 - mz)c(z)} = \frac{z^{k}}{z^{k} + (1 - 2z)c(z)}
\end{equation*}
The expected waiting time then follows immediately from the derivative at $z = 1/2$:
\begin{equation*}
    \left.z \frac{\dd}{\dd z} T(z)\right|_{z = 1/2} = 2^{k} c(1/2)
\end{equation*}

\section{Tree structures}
% Note I.39.
\noteblock{I.39}{Unary–binary trees and Motzkin numbers.}

Here the characteristic polynomial is $\phi(u) = 1 + u + u^{2}$.
To extract the number of trees with $n$ nodes, $M_{n} = [z^{n}] M(z)$, we apply Lagrange inversion:
\begin{equation*}
    \begin{aligned}
        M_{n}
        &= \frac{1}{n} [u^{n - 1}] \phi(u)^{n} \\
        &= \frac{1}{n} [u^{n - 1}] (1 + u + u^{2})^{n} \\
        &= \frac{1}{n} [u^{n - 1}] \sum_{j} \binom{n}{j}(u + u^{2})^{j} \\
        &= \frac{1}{n} \sum_{j} \binom{n}{j} [u^{n - 1}] u^{j}(1 + u)^{j} \\
        &= \frac{1}{n} \sum_{j} \binom{n}{j} [u^{n - j - 1}] \sum_{l} \binom{j}{l} u^{l} \\
        &= \frac{1}{n} \sum_{j} \binom{n}{j} \binom{j}{n - j - 1} \\
        &= \frac{1}{n} \sum_{k} \binom{n}{n - k} \binom{n - k}{n - (n - k) - 1} \\
        &= \frac{1}{n} \sum_{k} \binom{n}{k} \binom{n - k}{k - 1}
    \end{aligned}
\end{equation*}

% Note I.41.
\noteblock{I.41}{The Lagrangean form of Schröder’s GF.}

What remains is to prove the last equality:
\begin{equation*}
    \frac{(-1)^{n - 1}}{n} \sum_{k}(-2)^{k} \binom{n}{k + 1} \binom{n + k - 1}{k} =
        \frac{1}{n} \sum_{k = 0}^{n - 2} \binom{2n - k - 2}{n - 1} \binom{n - 2}{k}
\end{equation*}
Multiplying both sides by $n$ simplifies the notation:
\begin{equation*}
    (-1)^{n - 1} \sum_{k = 0}^{n - 1}(-2)^{k} \binom{n}{k + 1} \binom{n + k - 1}{k} = \sum_{k = 0}^{n - 2} \binom{2n - k - 2}{n - 1} \binom{n - 2}{k}
\end{equation*}

\stephead{1}{Re-index the left sum.}
Using $\binom{n}{k + 1} = \binom{n}{n - 1 - k}$ and $\binom{n + k - 1}{k} = \binom{n + k - 1}{n - 1}$, set $j = n - 1 - k$. Then
\begin{equation*}
    L = (-1)^{n - 1} \sum_{j = 0}^{n - 1}(-2)^{n - 1 - j} \binom{n}{j} \binom{2n - 2 - j}{n - 1} =
        2^{n - 1} \sum_{j = 0}^{n - 1} \left(-\frac{1}{2} \right)^{j} \binom{n}{j} \binom{2n - 2 - j}{n - 1}
\end{equation*}

\stephead{2}{Express as a coefficient extraction.}
Since $\binom{2n - 2 - j}{n - 1} = [z^{n - 1}] (1 + z)^{2n - 2 - j}$, coefficient extraction gives:
\begin{equation*}
    \begin{aligned}
        L
        &= 2^{n - 1} \sum_{j = 0}^{n - 1} \left(-\frac{1}{2} \right)^{j} \binom{n}{j} [z^{n - 1}] (1 + z)^{2n - 2 - j} \\
        &= 2^{n - 1}[z^{n - 1}] \sum_{j = 0}^{n - 1} \left(-\frac{1}{2} \right)^{j} \binom{n}{j}(1 + z)^{2n - 2 - j} \\
        &= 2^{n - 1}[z^{n - 1}] (1 + z)^{2n - 2} \sum_{j = 0}^{n} \binom{n}{j} \left(-\frac{1}{2(1 + z)} \right)^{j} \\
        &= 2^{n - 1}[z^{n - 1}] (1 + z)^{2n - 2}\left(1 - \frac{1}{2(1 + z)} \right)^{n} \\
        &= 2^{n - 1}[z^{n - 1}] (1 + z)^{2n - 2} \frac{(1 + 2z)^{n}}{2^{n}(1 + z)^{n}} \\
        &= \frac{1}{2} [z^{n - 1}] (1 + z)^{n - 2}(1 + 2z)^{n}
    \end{aligned}
\end{equation*}

\stephead{3}{Reverse the polynomial.}
Let $P(z) = (1 + z)^{n - 2}(1 + 2z)^{n}$, which has degree $2n - 2$.
The coefficient of $z^{n - 1}$ in $P(z)$ is the same as in the reversed polynomial $z^{2n - 2}P(1/z)$:
\begin{equation*}
    z^{2n - 2}P(1/z) = z^{2n - 2}\left(1 + \frac{1}{z} \right)^{n - 2}\left(1 + \frac{2}{z} \right)^{n} = (z + 1)^{n - 2}(z + 2)^{n}
\end{equation*}
Hence,
\begin{equation*}
    L = \frac{1}{2} [z^{n - 1}] (1 + z)^{n - 2}(2 + z)^{n}
\end{equation*}

\stephead{4}{Express the right sum as a coefficient.}
Likewise, $\binom{2n - k - 2}{n - 1} = [z^{n - 1}] (1 + z)^{2n - k - 2}$, so
\begin{equation*}
    \begin{aligned}
        R
        &= \sum_{k = 0}^{n - 2} \binom{2n - k - 2}{n - 1} \binom{n - 2}{k} \\
        &= \sum_{k = 0}^{n - 2} [z^{n - 1}] (1 + z)^{2n - k - 2} \binom{n - 2}{k} \\
        &= [z^{n - 1}] \sum_{k = 0}^{n - 2}(1 + z)^{2n - k - 2} \binom{n - 2}{k} \\
        &= [z^{n - 1}] (1 + z)^{2n - 2} \sum_{k = 0}^{n - 2} \binom{n - 2}{k}(1 + z)^{-k} \\
        &= [z^{n - 1}] (1 + z)^{2n - 2}\left(1 + \frac{1}{1 + z} \right)^{n - 2} \\
        &= [z^{n - 1}] (1 + z)^{2n - 2} \frac{(2 + z)^{n - 2}}{(1 + z)^{n - 2}} \\
        &= [z^{n - 1}] (1 + z)^{n}(2 + z)^{n - 2}
    \end{aligned}
\end{equation*}

\stephead{5}{Show equality of the two coefficient expressions.}
So we are down to showing that $\frac{1}{2}(1 + z)^{n - 2}(2 + z)^{n}$ and $(1 + z)^{n}(2 + z)^{n - 2}$ have the same coefficient of $z^{n - 1}$. Their difference is:
\begin{equation*}
    \begin{aligned}
        \Delta
        &= \frac{1}{2} [z^{n - 1}] (1 + z)^{n - 2}(2 + z)^{n} - [z^{n - 1}] (1 + z)^{n}(2 + z)^{n - 2} \\
        &= \frac{1}{2} [z^{n - 1}] (1 + z)^{n - 2}(2 + z)^{n - 2}[(2 + z)^{2} - 2(1 + z)^{2}] \\
        &= \frac{1}{2} [z^{n - 1}] (1 + z)^{n - 2}(2 + z)^{n - 2}(2 - z^{2})
    \end{aligned}
\end{equation*}
Thus it suffices to prove $[z^{n - 1}] f(z) = 0$, where
\begin{equation*}
    f(z) = (1 + z)^{n - 2}(2 + z)^{n - 2}(2 - z^{2})
\end{equation*}

\stephead{6}{A symmetry argument.}
$f(z)$ has degree $2n - 2$. Now compute $f(2 / z)$:
\begin{equation*}
    \begin{aligned}
        f(2/z)
        &= (1 + 2/z)^{n - 2}(2 + 2/z)^{n - 2}\left(2 - (2/z)^{2}\right) \\
        &= \left(\frac{2 + z}{z} \right)^{n - 2}\left(\frac{2(1 + z)}{z} \right)^{n - 2}\left(2 - \frac{4}{z^{2}}\right) \\
        &= \frac{(2 + z)^{n - 2}(1 + z)^{n - 2}2^{n - 2}}{z^{2n - 4}} \cdot \frac{-2(2 - z^{2})}{z^{2}} \\
        &= -2^{n - 1} \frac{f(z)}{z^{2n - 2}}
    \end{aligned}
\end{equation*}
Hence, $z^{2n - 2}f(2/z) = -2^{n - 1}f(z)$. Writing $f(z) = \sum_{k = 0}^{2n - 2} a_{k}z^{k}$ and comparing the middle coefficient gives:
$a_{n - 1}2^{n - 1} = -2^{n - 1}a_{n - 1}$, so necessarily $a_{n - 1} = 0$.

% Note I.42.
\noteblock{I.42}{Faster determination of Schröder numbers.}
\begin{equation*}
    S(z) = \frac{1}{4}(1 + z - \sqrt{1 - 6z + z^{2}})
\end{equation*}
The general template is
\begin{equation*}
    p(z)S(z) + q(z) = \sqrt[k]{r(z)}
\end{equation*}
with $p(z)$, $q(z)$, and $r(z)$ depending on $z$. Differentiate both sides:
\begin{equation*}
    p'(z)S(z) + p(z)S'(z) + q'(z) = \frac{1}{k}{r(z)}^{\frac{1}{k}-1}r'(z)
\end{equation*}
Then multiply through by $r(z)$:
\begin{equation*}
    p'(z)r(z)S(z) + p(z)r(z)S'(z) + q'(z)r(z) = \frac{1}{k}(p(z)S(z) + q(z))r'(z)
\end{equation*}
After rearranging,
\begin{equation*}
    p(z)r(z)S'(z) + (p'(z)r(z) - \frac{1}{k} p(z)r'(z))S(z) = \frac{1}{k} q(z)r'(z) - q'(z)r(z)
\end{equation*}
Now plug in $p(z) = -4$, $q(z) = 1 + z$, $k = 2$, and $r(z) = 1 - 6z + z^{2}$. A short simplification yields:
\begin{equation*}
    (z^{2} - 6z + 1)S'(z) + (3 - z)S(z) = 1 - z
\end{equation*}
from which the recurrence follows.

% Note I.43.
\noteblock{I.43}{Fast determination of the Cayley–Pólya numbers.}
\begin{equation*}
    H(z) = z \exp\left(H(z) + \frac{1}{2} H(z^{2}) + \cdots\right)
\end{equation*}

\stephead{1}{Take the natural logarithm.}
The logarithm brings the exponent down to a more manageable level:
\begin{equation*}
    \ln H(z) = \ln z + \sum_{k = 1}^{\infty} \frac{1}{k} H(z^{k})
\end{equation*}

\stephead{2}{Differentiate with respect to $z$.}
Differentiating gives:
\begin{equation*}
    \frac{H'(z)}{H(z)} = \frac{1}{z} + \sum_{k = 1}^{\infty} \frac{1}{k} H'(z^{k}) \cdot kz^{k - 1} = \frac{1}{z} + \sum_{k = 1}^{\infty} z^{k - 1}H'(z^{k})
\end{equation*}

\stephead{3}{Multiply by $zH(z)$.}
Multiplying by $zH(z)$ clears the fractions and lines up the powers of $z$:
\begin{equation*}
    zH'(z) = H(z)\left(1 + \sum_{k = 1}^{\infty} z^{k}H'(z^{k})\right)
\end{equation*}

\stephead{4}{Substitute power series.}
Now expand $H(z)$ as a formal power series $\sum_{n = 1}^{\infty} H_{n}z^{n}$. Then
\begin{itemize}
    \item $zH'(z) = \sum_{n = 1}^{\infty} n H_{n} z^{n}$
    \item $z^{k}H'(z^{k}) = \sum_{m = 1}^{\infty} m H_{m} z^{mk}$
\end{itemize}
Substituting these series into the equation,
\begin{equation*}
    \sum_{n = 1}^{\infty} n H_{n} z^{n} = H(z) + H(z)\sum_{k = 1}^{\infty} \sum_{m = 1}^{\infty} m H_{m} z^{mk}
\end{equation*}

\stephead{5}{The recurrence relation.}
Extracting the coefficient of $z^{n}$ leads to the recurrence. The double sum is a divisor-type convolution, and the standard form is
\begin{equation*}
    (n - 1)H_{n} = \sum_{l = 1}^{n - 1} H_{n - l}\left(\sum_{d|l} d H_{d}\right)
\end{equation*}
or, solving explicitly for $H_{n}$,
\begin{equation*}
    H_{n}
    = \frac{1}{n - 1} \sum_{l = 1}^{n - 1} H_{n - l}\left(\sum_{d|l} d H_{d}\right)
    = \frac{1}{n - 1} \sum_{l = 1}^{n - 1} a_{l}H_{n - l}
\end{equation*}
To see the recurrence at work, start from the base case $H_{1} = 1$, corresponding to the one-node tree.
\begin{itemize}
    \item To find $H_{2}$, we set $n = 2$. The summation runs for $l = 1$. The only divisor of $1$ is $1$.
        $a_{1} = \sum_{d|1} d H_{d} = 1 \cdot H_{1} = 1 \implies
        H_{2} = \frac{1}{2 - 1}(H_{2 - 1} \cdot a_{1}) = 1$.
    \item For $H_{3}$, the summation runs for $l = 1$ and $l = 2$. The divisors of $2$ are $1$ and $2$.
        $a_{2} = \sum_{d|2} d H_{d} = 1 \cdot H_{1} + 2 \cdot H_{2} = 3 \implies
        H_{3} = \frac{1}{3 - 1}(H_{3 - 1} \cdot a_{1} + H_{3 - 2} \cdot a_{2}) = 2$.
    \item For $H_{4}$, the summation runs for $l = 1, 2, 3$. The divisors of $3$ are $1$ and $3$.
        $a_{3} = \sum_{d|3} d H_{d} = 1 \cdot H_{1} + 3 \cdot H_{3} = 7 \implies
        H_{4} = \frac{1}{4 - 1}(H_{4 - 1} \cdot a_{1} + H_{4 - 2} \cdot a_{2} + H_{4 - 3} \cdot a_{3}) = 4$.
\end{itemize}

% Note I.50.
\noteblock{I.50}{Excursions, bridges, and meanders.}

A meander ending at altitude $k$ can be read as $k$ up-steps that create new floors, with Dyck excursions $\mathcal{D}$ inserted between consecutive floors and after the last one.
\begin{equation*}
    \mathcal{M}
        = \mathcal{D}
            + (\mathcal{D} \times \mathcal{Z} \times \mathcal{D})
            + (\mathcal{D} \times \mathcal{Z} \times \mathcal{D} \times \mathcal{Z} \times \mathcal{D})
            + \cdots
        = \mathcal{D} \times \sum_{k = 0}^{\infty}(\mathcal{Z} \times \mathcal{D})^{k}
\end{equation*}
A bridge is a sequence of primitive excursions. Each primitive excursion leaves 0, avoids returning in the middle, and comes back only at the end.
\begin{itemize}
    \item \textbf{Going Above:} Starts with +1, ends with -1, with a Dyck path in between: $\mathcal{Z} \times \mathcal{D} \times \mathcal{Z}$
    \item \textbf{Going Down:} Starts with -1, ends with 1, with a mirrored Dyck path in between: $\mathcal{Z} \times \mathcal{D} \times \mathcal{Z}$
\end{itemize}
\begin{equation*}
    \mathcal{B}
        = \mathrm{S}\textsc{eq}(\mathcal{Z} \times \mathcal{D} \times \mathcal{Z} + \mathcal{Z} \times \mathcal{D} \times \mathcal{Z})
        = \mathrm{S}\textsc{eq}(2 \times \mathcal{Z}^{2} \times \mathcal{D})
\end{equation*}
To prove $M_{n} = \binom{n}{\lfloor n/2 \rfloor}$, let $N(n, k)$ be the number of unrestricted paths of length $n$ ending at altitude $k$.
The reflection principle does the main work: paths that ever hit $-1$ and end at $k \geq 0$ are in bijection with unrestricted paths ending at $-k - 2$.
So the number of meanders ending at $k$ is $N(n, k) - N(n, k + 2)$, and summing over $k \geq 0$ gives the telescoping identity
\begin{equation*}
    M_{n} = \sum_{k \geq 0} [N(n, k) - N(n, k + 2)]
\end{equation*}
\begin{itemize}
    \item $n$ is even ($n = 2m$)
        $M_{n} = N(n, 0) = N(2m, 0) = \binom{2m}{\frac{2m + 0}{2}} = \binom{2m}{m} = \binom{n}{\lfloor n/2 \rfloor}$
    \item $n$ is odd ($n = 2m + 1$)
        $M_{n} = N(n, 1) = N(2m + 1, 1) = \binom{2m + 1}{\frac{2m + 1 + 1}{2}} = \binom{2m + 1}{m} = \binom{n}{\lfloor n/2 \rfloor}$
\end{itemize}

% Example I.17.
\exampleblock{I.17}{The complexity of boolean functions.}

We claim that
\begin{equation*}
    D_{n} = O(16^{n}m^{n}n^{-3/2})
\end{equation*}
implies
\begin{equation*}
    \sum_{j = 0}^{v} D_{j} = O(16^{v}m^{v}v^{-3/2})
\end{equation*}
The clean way is to prove a slightly more general statement first. \\
Let $c > 1$ and $\alpha > 0$. Define $a_{n} = c^{n}n^{-\alpha}$ for $n \geq 1$ and let $a_{0} > 0$ be a constant.
If $D_{n} = O(a_{n})$, then $\sum_{j = 0}^{v} D_{j} = O(c^{v}v^{-\alpha})$ as $v \to \infty$. \\

Proof. Since $D_{n} = O(a_{n})$, there exist constants $M > 0$ and $N \geq 1$ such that $D_{n} \leq Mc^{n}n^{-\alpha}$ for all $n \geq N$.
For $v \geq N$ we split:
\begin{equation*}
    \sum_{j = 0}^{v} D_{j} \leq \underbrace{\sum_{j = 0}^{N - 1} D_{j}}_{C_{0}} + M \sum_{j = N}^{v} c^{j}j^{-\alpha}
\end{equation*}
So the problem reduces to controlling
\begin{equation*}
    S(v) = \sum_{j = N}^{v} c^{j}j^{-\alpha}.
\end{equation*}
Split the sum at $j = \lfloor v/2 \rfloor$:
\begin{equation*}
    S(v) = \underbrace{\sum_{j = N}^{\lfloor v/2 \rfloor} c^{j}j^{-\alpha}}_{\text{early}}
        + \underbrace{\sum_{j = \lfloor v/2 \rfloor + 1}^{v} c^{j}j^{-\alpha}}_{\text{tail}}
\end{equation*}
\textbf{Early terms.} For $j \leq v/2$ we have $c^{j} \leq c^{v/2}$. Also since $\alpha > 0$, the factor $j^{-\alpha}$ is decreasing;
its maximum is at $j = N$, so $j^{-\alpha} \leq N^{-\alpha}$. Thus
\begin{equation*}
    \sum_{j = N}^{\lfloor v/2 \rfloor} c^{j}j^{-\alpha}
        \leq N^{-\alpha} \sum_{j = N}^{\lfloor v/2 \rfloor} c^{j}
        \leq N^{-\alpha} \frac{c^{v/2 + 1}}{c - 1} = O(c^{v/2})
\end{equation*}
Because $c > 1$, $c^{v/2}$ is exponentially smaller than $c^{v}$, so this part is negligible on the relevant scale. \\
\textbf{Tail terms.} For $j \geq v/2$, $j^{-\alpha} \leq (v/2)^{-\alpha} = 2^{\alpha}v^{-\alpha}$. Consequently
\begin{equation*}
    \sum_{j = \lfloor v/2 \rfloor + 1}^{v} c^{j}j^{-\alpha}
        \leq 2^{\alpha}v^{-\alpha} \sum_{j = \lfloor v/2 \rfloor + 1}^{v} c^{j}
        \leq 2^{\alpha}v^{-\alpha} \frac{c^{v + 1}}{c - 1} = O(c^{v}v^{-\alpha})
\end{equation*}
Putting the two pieces together, we get $S(v) = O(c^{v}v^{-\alpha})$, and hence
\begin{equation*}
    \sum_{j = 0}^{v} D_{j} = C_{0} + MS(v) = O(1) + O(c^{v}v^{-\alpha}) = O(c^{v}v^{-\alpha})
\end{equation*}

With the general estimate in hand, return to the original problem. We now show
\begin{equation*}
    \sum_{j = 0}^{v} D_{j} = O(16^{v}m^{v}v^{-3/2}) = o(\lVert \Omega \rVert)
\end{equation*}
by taking for $v$ the value
\begin{equation*}
    v = \frac{2^{m}}{4 + \log_{2}m}
\end{equation*}
Observe that
\begin{equation*}
    v\log_{2}(16m) = \frac{2^{m}}{4 + \log_{2}m}(4 + \log_{2}m) = 2^{m} \implies 16^{v}m^{v} = 2^{2^{m}}
\end{equation*}
So the bound simplifies to
\begin{equation*}
    \sum_{j = 0}^{v} D_{j} = O(2^{2^{m}}v^{-3/2})
\end{equation*}
Dividing by $\lVert \Omega \rVert$:
\begin{equation*}
    \frac{\sum_{j = 0}^{v} D_{j}}{\lVert \Omega \rVert} = O(v^{-3/2})
\end{equation*}
As $m \to \infty$, we also have $v \to \infty$, hence $v^{-3/2} \to 0$. Therefore,
\begin{equation*}
    \frac{\sum_{j = 0}^{v} D_{j}}{\lVert \Omega \rVert} \to 0
\end{equation*}
which is exactly $\sum_{j = 0}^{v} D_{j} = o(\lVert \Omega \rVert)$.

\section{Additional constructions}
% Note I.62.
\noteblock{I.62}{Sequences without repeated components.}

A sequence without repeated components is an ordered arrangement of distinct objects drawn from a base class $\mathcal{B}$.
The natural route is to build the set first and recover the ordering afterward. The generating function for the power-set construction $\mathrm{PS}\textsc{et}(\mathcal{B})$ is
\begin{equation*}
    A(z, u)
        = \prod_{\beta \in \mathcal{B}}\left(1 + uz^{|\beta|}\right)
        = \exp\left(\sum_{j \geq 1}(-1)^{j - 1} \frac{u^{j}}{j}B(z^{j})\right)
\end{equation*}
Here $[u^{k}] A(z, u) = f^{\langle k \rangle}(z)$ is the generating function for sets with exactly $k$ distinct components.
The no-repeat condition is already built in, but order is not. To turn sets into sequences, we multiply by $k!$ to account for all permutations.
Using the Eulerian integral representation $k! = \int_{0}^{\infty} e^{-u}u^{k}\,\dd u$, we can write the generating function for distinct sequences,
$\mathcal{S}_{\text{distinct}}(z)$, as follows:
\begin{equation*}
    \begin{aligned}
        \mathcal{S}_{\text{distinct}}(z)
            &= \sum_{k \geq 0}(\text{Sets of size k}) \cdot k! \\
            &= \sum_{k \geq 0} f^{\langle k \rangle}(z) \cdot k! \\
            &= \sum_{k \geq 0} f^{\langle k \rangle}(z) \cdot \int_{0}^{\infty} e^{-u}u^{k}\,\dd u \\
            &= \sum_{k \geq 0} \int_{0}^{\infty} f^{\langle k \rangle}(z)e^{-u}u^{k}\,\dd u \\
            &= \int_{0}^{\infty} \sum_{k \geq 0} f^{\langle k \rangle}(z)e^{-u}u^{k}\,\dd u \\
            &= \int_{0}^{\infty} \sum_{k \geq 0} f^{\langle k \rangle}(z)u^{k} \cdot e^{-u}\,\dd u \\
            &= \int_{0}^{\infty} A(z, u) \cdot e^{-u}\,\dd u \\
            &= \int_{0}^{\infty} \exp\left(\sum_{j \geq 1}(-1)^{j - 1} \frac{u^{j}}{j}B(z^{j})\right)e^{-u}\,\dd u
    \end{aligned}
\end{equation*}

\section{Perspective}
